% sections/02_institutions.tex -- Section 2 (W3). Plan: BLUEPRINT 2.2 and 5.3; Overrides O4, O5, O8, O9.
% Provenance convention (Override O5): "% src: <file>:<line>" follows every sentence that carries a number.
% Abbreviations in src comments: 12:NN = frot/sub2025/12_app_quant.tex; 7:NN = 7_policy_reforms.tex; 5:NN = 5_quantitative_model.tex;
% 1:NN = 1_intro.tex; T: = paper/tables/; v6: = frot/current/.
\section{Social Security Rules, the Financing Problem, and the Reform Menu}\label{sec:institutions}

\subsection{Rules faced by the 1920--1940 cohorts}\label{subsec:rules}

Individuals born between 1920 and 1940 could first claim Old-Age, Survivors, and Disability Insurance (OASDI) retirement benefits at the early retirement age (ERA) of 62 and faced, on average over birth years, a normal retirement age (NRA) of 65. % src: 12:10; T:parameters_ss.tex:46-47
The benefit is the primary insurance amount (PIA) multiplied by a claiming-age factor that falls by 6.67 percentage points for each year of claiming before the NRA and rises by the delayed retirement credit (DRC) for each year of claiming after it, up to age 70. The DRC ranged from 3 to 6 percent per year by year of birth, and the paper uses the cohort average of 4.5 percent. % src: 7:90; 12:10
\autoref{tab:ssdrc} reports the schedule: 80 percent of the PIA at age 62, 100 percent at 65, and 122.5 percent at 70. % src: T:SS_DRC.tex:10

\begin{table}[!t]
\centering
\caption{Social Security Benefit as a Percentage of the Primary Insurance Amount by Claiming Age, 1920--1940 Cohorts (NRA 65, DRC 4.5 Percent)}
\label{tab:ssdrc}
\vspace{6pt}
\footnotesize
\begin{threeparttable}
    \begin{tabular}{l|l|l l l l l l l l l} % <-- Alignments: 1st column left, 2nd middle and 3rd right, with vertical lines in between
 \hline\hline
%  & & & & {}\\
  &  & \multicolumn{9}{c}{Benefit, as a percentage of PIA, payable at age} \\ [0.5ex]
 \hline
%&& & & & {}\\
NRA & DRC	& 62&63&64&65&66&67&68&69&70\\[0.5ex]
\hline
%&& & & & {}\\
65  & 4.5  & 80 & 86.67&93.33&	100	&104.5&	109&	113.5 & 118&122.5\\
%66  & 8 &  75	 &80	 &86.67 &	93.33 &	100 &	108 &116 &124 &	132\\[0.5ex] 
%&67 & & 8& &70 &	75	&80&	86.67&	93.33&	100&108&116 &124 \\[1ex] 
 \hline\hline
    \end{tabular}

\begin{tablenotes}[flushleft]\scriptsize
\item[] \textit{Notes:} Benefit payable at ages 62--70 as a percentage of the PIA, from the rules of the Social Security Administration; the reduction before the NRA is 6.67 percentage points per year and the delayed retirement credit after the NRA is 4.5 percentage points per year. NRA, normal retirement age; DRC, delayed retirement credit, in percent per year. Source: Social Security Administration.
% src: T:SS_DRC.tex:10; 7:90
\end{tablenotes}
\end{threeparttable}
\end{table}

The PIA applies factors of 0.9, 0.32, and 0.15 to average indexed monthly earnings (AIME), the average of the highest 35 years of indexed earnings, below \$3,720, between \$3,720 and \$22,392, and above \$22,392, with AIME capped at \$43,800 (1987 dollars, as are the other amounts in this subsection). % src: 5:70; 5:102-104; T:parameters_ss.tex:43-45; 12:12
Under the retirement earnings test (RET), a beneficiary loses one dollar of benefits for every two dollars of earnings above \$6,000 before the NRA and one for every three dollars above \$8,186 from the NRA to age 69, with no test from age 70; withheld benefits are credited to later benefits \citep{french2011effects}. % src: T:parameters_ss.tex:50-52; 12:12; 5:140; 5:148-149
Disability insurance (DI) pays a PIA-based benefit, without the early-claiming reduction, to a worker with a disability who is older than 50, younger than the NRA, and earns at most \$3,600 a year. % src: 12:46-49; T:parameters_ss.tex:42
A payroll tax of 3 to 7.65 percent on earnings up to \$14,459--\$64,974 finances OASDI and Hospital Insurance. % src: T:parameters_ss.tex:37-38; 12:89
\autoref{tab:ssparams} collects these parameters; \autoref{sec:model} states how each enters the individual's problem.

\subsection{A higher normal retirement age or a steeper actuarial adjustment}\label{subsec:twoways}

The schedule of \autoref{tab:ssdrc} is the factor $\psi(a^c)$ that multiplies the PIA in \eqref{eq:adjust}. A reform can change it in two ways.

The first way moves the NRA. Reforms AN1--AN3 raise it to 66, 67, and 68, and the whole schedule moves one year later for each year added: the factor at age 62 falls from 80 to 75, 70, and 65 percent of the PIA, and the factor at age 70 from 122.5 to 118, 113.5, and 109. % src: T:DRC_visualization_4.tex:37-40
Every claimant receives a lower benefit at a given claiming age, whether or not the claiming age changes. % src: 7:90
This is the instrument of the 1983 amendments, which raise the NRA to 67 for cohorts born in 1960 and later, and of the proposal, analyzed by \citet{cbo2024socialsecurity}, to raise it to 70 for cohorts born in 1981 and later. % src: 1:44 (footnote)

The second way leaves the factor at the NRA at 100 percent and changes the slope of the schedule around it. Reforms AD1--AD5 raise the DRC from the cohort's 4.5 percent to 8, 8.5, 9, 9.5, and 10 percent per year and leave the reduction before the NRA unchanged, so the factor at age 70 rises from 122.5 to 140 under AD1 and 150 under AD5 while the factors at ages 62--65 are those of the status quo. % src: T:DRC_visualization_4.tex:37,43-47
Reforms AL1--AL9 set the early-claiming reduction and the DRC equal in absolute value, at $\pm 7$ to $\pm 15$ percent per year in steps of one percentage point; the factor at age 62 falls from 79 under AL1 to 55 under AL9, and the factor at age 70 rises from 135 to 175. % src: T:DRC_visualization_4.tex:50,58
Under both the AD and the AL reforms a claimant at the NRA receives the status-quo benefit; what changes is the price of claiming earlier or later. % src: 7:92
AN1 and AD1 applied together, an NRA of 66 and a DRC of 8 percent, are the rules that the cohort born in the 1950s faced in \citet{Yu2024}; here each is applied alone to the cohort that faced an NRA of 65 and a DRC of 4.5 percent, and the combination was not simulated. % src: frot/yu_jpube/Yu_Zhixiu_retirement_NBER.tex:387-388; T:parameters_ss.tex:47-48
% TODO-AUTHORS: Open question 18 -- verify that the published JPubE table prints NRA 66 / DRC 8 for the 1950s cohort; the sentence is written to hold for the NBER version (Override O9).

\subsection{The reform menu}\label{subsec:menu}

\autoref{tab:menu} lists every reform evaluated in \autoref{sec:claiming} to \autoref{sec:other}.\footnote{There are 66 reform rows in seven families; L9 and P3.9 are the same policy, and D3 and K3 are the same experiment.} % src: rows of T:welfare_age_dependence_mvpfv5.tex:108-126, T:welfare_progressivity_mvpfv7.tex:115-158, T:welfare_medicare_mvpfv8.tex:113,116, T:welfare_tax_mvpf_v8.tex:114-118, T:welfare_savings_mvpf_v3.tex:120-123; wording fixed in BLUEPRINT Decision 9
The benefit-formula reforms raise the first, second, or third PIA factor in steps of one percent of its initial value (P1.1--P1.9, P2.0--P2.9, P3.0--P3.9) or all three factors proportionally (L0--L9); L9 and P3.9 set the same factors. % src: 12:475; T:welfare_progressivity_mvpfv7.tex:115-158
Two Medicare reforms raise the share of total medical spending covered after age 65 by 20 percentage points, for everyone (Part A) or for nonparticipants only (Part B), because employer-provided plans already cover physician and outpatient services for most workers. % src: 12:536; 7:198; T:welfare_medicare_mvpfv8.tex:113,116
% TODO-AUTHORS: Open question 5 -- confirm the Medicare unit (percentage points of total spending; table label "+ 20 pp"; May 2025 text "by 20%").
Income-tax reforms I1--I5 lower income taxes after age 62 so that after-tax income rises by 1 to 5 percent. % src: 7:349; T:welfare_tax_mvpf_v8.tex:114-118
% TODO-AUTHORS: Open question 6 -- 12:557 implements lambda_t + x(t-61), which grows with age; confirm which rule was run.
Saving-subsidy reforms K1--K3 raise the net return on assets by 0.2, 0.4, and 0.8 percentage points; they stand for an expansion of the tax preferences that retirement saving already receives through individual retirement accounts and 401(k) plans, and K3, which brings the net return of 3.2 percent to the gross return of 4 percent, is the experiment reported as D3 in \autoref{sec:elimination}. % src: 12:560; 5:390; 7:423 (IRAs and 401(k) plans)
\autoref{app:reforms} gives the implementation of each reform.

% tables/new_reform_menu.tex -- Table 2 (tab:menu), written by W3. Text-only table of the reform families (BLUEPRINT 2.2).
% Every entry below is taken from the following sources (no number is computed here):
%   AN1-AN3, AD1-AD5, AL1-AL9 definitions:     frot/sub2025/12_app_quant.tex:487 (verbatim in the comment block after the notes; not rendered); grids as printed in tables/DRC_visualization_4.tex:38-40, 43-47, 50-58
%   status-quo NRA 65 / DRC 4.5 percent:        tables/parameters_ss.tex:47-48; 12_app_quant.tex:10
%   P1.1-P1.9, P2.0-P2.9, P3.0-P3.9, L0-L9:      frot/sub2025/12_app_quant.tex:475 (verbatim in the comment block after the notes; not rendered); bracket triples as printed in
%                                               tables/welfare_progressivity_mvpfv7.tex:115 (P1.1: 0.91,0.32,0.15), :123 (P1.9: 0.99,0.32,0.15), :126 (P2.0: 0.99,0.3232,0.15),
%                                               :135 (P2.9: 0.99,0.352,0.15), :137 (P3.0: 0.99,0.352,0.1515), :146 (P3.9: 0.99,0.352,0.165), :149 (L0: 0.909,0.3232,0.1515), :158 (L9: 0.99,0.352,0.165)
%   L9 = P3.9:                                  tables/welfare_progressivity_mvpfv7.tex:146 versus :158 (same triple, same outcomes)
%   Medicare Part A / Part B:                   frot/sub2025/12_app_quant.tex:536 (verbatim in the comment block after the notes; not rendered); 7_policy_reforms.tex:198; row labels "+ 20 pp" in tables/welfare_medicare_mvpfv8.tex:113,116
%   I1-I5:                                      frot/sub2025/12_app_quant.tex:556-557 (verbatim in the comment block after the notes; not rendered); row labels in tables/welfare_tax_mvpf_v8.tex:114-118; 7_policy_reforms.tex:349
%   K1-K3:                                      frot/sub2025/12_app_quant.tex:560 (verbatim in the comment block after the notes; not rendered); row labels in tables/welfare_savings_mvpf_v3.tex:120-123
%   D3 = K3:                                    tables/welfare_decomposition_gradient.tex:33 (3.60) versus tables/welfare_savings_mvpf_v3.tex:123 (3.6); 5_quantitative_model.tex:390 (3.2 -> 4 percent); D03 section 6 item 9
%TODO-AUTHORS: Open question 5 -- the Medicare rows say "percentage points" (table label "+ 20 pp"); the May 2025 text says "by 20%". Confirm x = 0.20 of total spending.
%TODO-AUTHORS: Open question 6 -- the I rows follow the table labels (after-tax income +1 to +5 percent); the appendix formula lambda_t + x(t-61) grows with age. Confirm which was run.
\begin{table}[!t]
\centering
\caption{The Reform Menu}
\label{tab:menu}
\vspace{6pt}
\scriptsize
\setlength{\tabcolsep}{4pt}
\renewcommand{\arraystretch}{1.05}
\begin{threeparttable}
\begin{tabular}{@{}>{\raggedright\arraybackslash}p{1.05in}>{\raggedright\arraybackslash}p{0.7in}>{\raggedright\arraybackslash}p{1.75in}>{\raggedright\arraybackslash}p{1.6in}>{\raggedright\arraybackslash}p{1.05in}@{}}
\toprule
Family & IDs & Instrument changed & Grid & Ages affected \\
\midrule
Normal retirement age & AN1--AN3 & NRA; the claiming schedule of \autoref{tab:ssdrc} moves one year later per step & NRA 66, 67, 68 (status quo 65) & Every claiming age 62--70 \\
\addlinespace
Delayed retirement credit & AD1--AD5 & DRC, the credit per year of claiming after the NRA; the reduction before the NRA is unchanged & 8, 8.5, 9, 9.5, 10 percent per year (status quo 4.5) & Claiming ages 66--70 \\
\addlinespace
Symmetric adjustment & AL1--AL9 & Early-claiming reduction and DRC set equal in absolute value & $\pm7$ to $\pm15$ percent per year, in 1 pp steps & Claiming ages 62--64 and 66--70 \\
\addlinespace
Benefit formula: progressivity & P1.1--P1.9 & First PIA factor & 0.91 to 0.99, 1 pp per step (status quo 0.90) & All beneficiaries, from the claiming age \\
 & P2.0--P2.9 & Second PIA factor, with the first factor at 0.99 & 0.3232 to 0.352, 0.32 pp per step (status quo 0.32) & All beneficiaries, from the claiming age \\
 & P3.0--P3.9 & Third PIA factor, with the first and second factors at 0.99 and 0.352 & 0.1515 to 0.165, 0.15 pp per step (status quo 0.15) & All beneficiaries, from the claiming age \\
\addlinespace
Benefit formula: level & L0--L9 & All three PIA factors & $+1$ percent of each factor per step, from (0.909, 0.3232, 0.1515) to (0.99, 0.352, 0.165); L9 coincides with P3.9\tnote{a} & All beneficiaries, from the claiming age \\
\addlinespace
Medicare & Part A & Share of total medical spending covered, everyone aged 65 and older & $+20$ percentage points & 65 and older \\
 & Part B & Share of total medical spending covered, nonparticipants aged 65 and older & $+20$ percentage points & 65 and older, nonparticipants \\
\addlinespace
Income tax & I1--I5 & Income tax after age 62; after-tax income $y_t-T_t(y_t)$ rises & $+1$ to $+5$ percent, in 1 pp steps & 62 and older \\
\addlinespace
Saving subsidy & K1--K3 & Net return on assets $r$ & $+0.2$, $+0.4$, $+0.8$ pp (status quo 3.2 percent); K3 coincides with D3\tnote{b} & All ages \\
\bottomrule
\end{tabular}
\begin{tablenotes}[flushleft]\scriptsize
\item[] \textit{Notes:} NRA, normal retirement age; DRC, delayed retirement credit; PIA, primary insurance amount; pp, percentage points. The status-quo values are those of the 1920--1940 cohorts (\autoref{tab:ssdrc} and \autoref{tab:ssparams}). The rule change behind each family is stated in Appendix~\ref{app:construction}; the Medicare rows are in percentage points of the coverage share of total medical spending.
% src: tables/parameters_ss.tex:46-48 (status quo); appendix/C_reforms.tex (app:construction); tables/welfare_medicare_mvpfv8.tex:113,116 ("+ 20 pp")
% ---- Provenance (not rendered): the May 2025 construction definitions, quoted verbatim, that the rows above and Appendix C.1 paraphrase.
% ---- Moved out of the rendered note on 2026-11-09 (verification pass): the note ran to 27 lines of \scriptsize, duplicated Appendix C.1,
% ---- and the Medicare quotation says "by 20%" where the rows, the text, and Appendix C say 20 percentage points (BLUEPRINT 2.2 do-not-say; Open question 5).
% 12_app_quant.tex:475: "Social Security progressivity reforms P1.1--P1.9 modify the first bracket of the primary insurance amount (PIA) formula by gradually increasing its replacement rate. Specifically, the initial rate of 0.90 is raised by one percentage point (pp) in each step, from 0.91 under P1.1 up to 0.99 under P1.9. Reforms P2.0--P2.9 and P3.0--P3.9 target the second and third brackets, respectively: P2.0--P2.9 increase the second bracket's replacement rate by 0.32 pp per step, while P3.0--P3.9 raise the third bracket's replacement rate by 0.15 pp per step. Each set of reforms implements changes in 1% increments relative to the initial bracket levels. In addition, Social Security level reforms L0--L9 uniformly increase the replacement rates across all PIA brackets by 1% in each step."
% 12_app_quant.tex:487: "Reforms that modify the NRA, denoted as AN1--3, increase the NRA to ages 66, 67, and 68, respectively. Reforms that adjust the DRC, referred to as AD1--5, increase the DRC to 8%, 8.5%, 9%, 9.5%, and 10%, respectively. Reforms that change the linear adjustment of retirement benefits (AL1--9) modify both the DRC and early retirement penalties, increasing both rates from 7% to 15%, with a 1 pp increase for each reform step."
% 12_app_quant.tex:536: "Medicare Part A reform increases coverage for all individuals over age 65 by 20%, reflecting expanded access to inpatient services; Medicare Part B reform increases coverage among nonparticipants by 20%, improving access to outpatient care for those not previously enrolled" (coverage rates for individuals aged 65 and older computed by health and working status from MEPS data; the model re-simulated with these rates is the baseline for the two reforms).
% 12_app_quant.tex:556-557: "Income tax reforms (I1--I5) reduce income taxes after age 62, leading to increases in after-tax income y-T(y)=lambda_t y^(1-tau_t)."
% 12_app_quant.tex:560: "Savings subsidy reforms K1--3 increase the net returns to savings by 0.20 pp, 0.4 pp, and 0.8 pp."
\item[a] L9 and P3.9 set the same three factors, (0.99, 0.352, 0.165), and are the same policy; \autoref{tab:progressivity_full} reports identical outcomes for the two rows.
% src: tables/welfare_progressivity_mvpfv7.tex:146 versus :158
\item[b] K3 raises the net return from 3.2 to 4 percent, the gross return; it is the experiment reported as D3 in \autoref{tab:elim}.
% src: 5_quantitative_model.tex:390; tables/welfare_decomposition_gradient.tex:33; tables/welfare_savings_mvpf_v3.tex:123
\end{tablenotes}
\end{threeparttable}
\end{table}


\subsection{The cohort perspective and the financing benchmark}\label{subsec:cohort}

Each reform is a permanent change in the rules faced by one cohort from age 30. Welfare is the ex ante average of expected utility at age 30 (\autoref{sec:accounting}). Fiscal flows are the cohort's own---benefits, Medicare payments, and transfers received, less payroll and income taxes paid---as per-capita net present values at age 62 at the net return of 3.2 percent; a reform changes the cohort's net position by the change in net government spending that the Net column of every reform table measures (\autoref{app:derivations}). % src: 12:611; 5:390
Fiscal savings are valued at face value: a dollar the cohort does not receive is available to taxpayers or beneficiaries outside it, on whom we place a unit welfare weight. Incidence within the cohort is valued by the ex ante average, not by the distribution of gains and losses across realized histories. Incidence across cohorts is not weighted here: the marginal value of public funds is a shadow price, the welfare delivered to the cohort per dollar of net spending on it \citep{hendren2020unified}, to which any intergenerational weight can be applied through the Net column. No transition is modeled, and nothing is said about the welfare of other cohorts.

The benchmark on the fiscal side is the long-run financing gap: the 2024 Trustees Report puts the 75-year actuarial deficit of OASDI at 3.50 percent of taxable payroll \citep{trustees2024}. % src: v6:6_quant...resubmission_v6.tex:546 (figure carried from the companion paper)
% Trustees figure verified against the 2024 OASDI Trustees Report by web search on 2026-11-09 (Override O4; AUTHOR_NOTES.md Open question 13; same check as sections/01_intro.tex, Section 1.1).
A reform contributes to closing it if its net fiscal effect on the cohort is negative. The Net column is in percent of the cohort's baseline net spending; expressing it in percent of taxable payroll would require the present value of the cohort's payroll, which the reported baseline aggregates do not include, so the comparison with the Trustees' figure is one of sign.
% TODO-AUTHORS: Open question 7 -- a payroll-unit conversion needs PV_62 of the cohort's taxable payroll from baseline output.
